\subsection{AoI-oriented UAV data collection}
The Age of Information (AoI) was introduced in \cite{kaul2012aoi} as a
destination-centric measure of data freshness, capturing the time elapsed
since the generation of the most recently received update. AoI has since
become a standard objective for UAV-assisted data collection, where a UAV
visits distributed sources and the design goal is to keep the collected
information fresh. In the swarm setting, \cite{eldeeb2023swarm} minimizes AoI
across multiple cooperating UAVs using multi-agent reinforcement learning,
demonstrating that coordinated trajectories can reduce peak AoI relative to
uncoordinated flight. That formulation, however, assumes flight endurance is
unconstrained: it carries no energy state and no recharging model, so the
depletion and replenishment dynamics that dominate long missions are absent.
Closest to our objective is \cite{base2024}, which jointly optimizes charging
station placement, UAV trajectory, and peak AoI, and which we adopt as our
point of departure. That work is nonetheless single-UAV: with one vehicle
there is no inter-UAV contention for a charger, no need for multi-port
capacity, and no swarm assignment, precisely the coupling that our design
targets.

\subsection{Energy-limited UAVs and charging}
A second line of work makes energy replenishment explicit. In \cite{wei2024},
multiple UAVs share a limited charging facility and the resulting charge
waiting time is folded into the AoI objective, then optimized by deep
reinforcement learning. The facility there is single and fixed in location,
which effectively reduces contention to a single-server queue and leaves
station placement out of scope. Laser-based (LBD) charging is studied in
\cite{sun2025laser}, where charging sites are at fixed locations and are
treated without queue contention, so neither where to place stations nor how
many vehicles may wait is optimized. A different tactic uses mobile chargers:
\cite{oubbati2024ugv} and \cite{zhao2025aoi} employ ground vehicles that
follow the UAVs to deliver energy on the move. Because the charger comes to
the UAV, this removes the fixed-station contention that is central to our
setting and does not address the placement of fixed, multi-port stations.
Finally, \cite{ev2022sharing} does model UAV charging stations as a contended
queue and provides a waiting-time analysis, but it optimizes neither AoI nor
trajectory nor station placement, so the queueing model is not coupled to any
freshness or mobility decision.

\subsection{Positioning and gap}
Read together, these works each supply a subset of the ingredients our problem
requires but never their combination. Swarm operation with AoI appears in
\cite{eldeeb2023swarm} yet without charging; placement with AoI appears in
\cite{base2024} yet only for a single UAV; charge waiting is captured in
\cite{wei2024} yet at a single fixed server without placement; and a contended
charging queue is analyzed in \cite{ev2022sharing} yet without AoI,
trajectory, or placement. No prior formulation simultaneously carries all four
of station placement, swarm operation, multi-server charging contention, and a
peak-AoI objective. Our work closes this gap by jointly optimizing station
placement, port capacity, swarm trajectory, and source-to-UAV assignment under
a charging model that is a contended, multi-server, and finite-source (closed)
queue, validated against discrete-event simulation. The finite-source aspect
is not a cosmetic refinement: because the number of UAVs competing for a
station is small and closed, modeling the queue as an open M/M/c system
over-predicts waiting time and thereby distorts both placement and capacity
decisions, an effect we quantify and correct with the finite-source model
detailed in Section~\ref{sec:model}. This coupled formulation also yields two
structural findings that a decoupled treatment cannot expose: a congestion
threshold, beyond which peak AoI is non-monotone in swarm size, and a
placement inversion, whereby coverage-optimal station placement becomes
AoI-suboptimal once contention is accounted for.
